\begin{resumen}

    La inspección de calidad en entornos industriales se enfrenta al desafío crítico de la escasez de muestras de defectos ("Small Data") y la ineficiencia de la supervisión manual. En esta investigación, hemos analizado las limitaciones de las redes convolucionales (CNN) nativas, las cuales presentaron resultados insuficientes en la detección de anomalías milimétricas. 
    
    Para resolver esta problemática, proponemos la implementación de Vision Transformers (\cite{dosovitskiy2021vit}) mediante la arquitectura DEIMv2, potenciada por un pre-entrenamiento auto-supervisado basado en el hito DINOv3~\cite{simeoni2025dinov3}. A diferencia de las soluciones locales tradicionales, nuestro enfoque utiliza mecanismos de atención global para capturar discontinuidades estructurales en los materiales.
    
    Los resultados experimentales validan la superioridad de la propuesta, alcanzando un mAP (\textit{Mean Average Precision}) de 0.785. Mediante una rigurosa validación científica en la Fase 3, hemos demostrado que esta mejora es intrínseca a la arquitectura del Vision Transformer y no a un incremento en la resolución de entrada. La robustez del sistema, confirmada bajo diversos umbrales de decisión, posiciona esta metodología como un estándar viable para la automatización de la calidad industrial.
    
    Este trabajo sienta las bases para el análisis detallado del contexto industrial y las soluciones tecnológicas que se presentan en el Capítulo 1.
    \end{resumen}
    
    \begin{abstract}
    \noindent Industrial quality inspection faces a critical challenge due to the scarcity of defect samples ("Small Data") and the inherent subjectivity of manual supervision. This research analyzes the limitations of native Convolutional Neural Networks (CNNs), which demonstrated insufficient performance in detecting millimeter-scale anomalies.
    
    To address this issue, we propose the implementation of Vision Transformers (\cite{dosovitskiy2021vit}) using the DEIMv2 architecture, enhanced by a self-supervised pre-training strategy based on the DINOv3 milestone~\cite{simeoni2025dinov3}. Unlike traditional local-based solutions, our approach leverages global attention mechanisms to identify structural discontinuities in industrial materials.
    
    The experimental results validate the superiority of the proposed method, achieving a mAP (Mean Average Precision) of 0.785. Through a rigorous scientific validation in Phase 3, we have proved that this performance gain is intrinsic to the Vision Transformer architecture rather than a mere increase in input resolution. The system's robustness, confirmed across various decision thresholds, establishes this methodology as a viable standard for industrial quality automation.
    
    This abstract provides a concise overview of the investigation, which is further detailed starting with the industrial context in Chapter 1.
\end{abstract}